% Options for packages loaded elsewhere
\PassOptionsToPackage{unicode}{hyperref}
\PassOptionsToPackage{hyphens}{url}
%
\documentclass[
  11pt,
  a4paper,
]{article}
\usepackage{amsmath,amssymb}
\usepackage{iftex}
\ifPDFTeX
  \usepackage[T1]{fontenc}
  \usepackage[utf8]{inputenc}
  \usepackage{textcomp} % provide euro and other symbols
\else % if luatex or xetex
  \usepackage{unicode-math} % this also loads fontspec
  \defaultfontfeatures{Scale=MatchLowercase}
  \defaultfontfeatures[\rmfamily]{Ligatures=TeX,Scale=1}
\fi
\usepackage{lmodern}
\ifPDFTeX\else
  % xetex/luatex font selection
  \setmainfont[]{Latin Modern Roman}
  \setmathfont[]{Latin Modern Math}
\fi
% Use upquote if available, for straight quotes in verbatim environments
\IfFileExists{upquote.sty}{\usepackage{upquote}}{}
\IfFileExists{microtype.sty}{% use microtype if available
  \usepackage[]{microtype}
  \UseMicrotypeSet[protrusion]{basicmath} % disable protrusion for tt fonts
}{}
\makeatletter
\@ifundefined{KOMAClassName}{% if non-KOMA class
  \IfFileExists{parskip.sty}{%
    \usepackage{parskip}
  }{% else
    \setlength{\parindent}{0pt}
    \setlength{\parskip}{6pt plus 2pt minus 1pt}}
}{% if KOMA class
  \KOMAoptions{parskip=half}}
\makeatother
\usepackage{xcolor}
\usepackage[margin=24mm]{geometry}
\usepackage{longtable,booktabs,array}
\usepackage{calc} % for calculating minipage widths
% Correct order of tables after \paragraph or \subparagraph
\usepackage{etoolbox}
\makeatletter
\patchcmd\longtable{\par}{\if@noskipsec\mbox{}\fi\par}{}{}
\makeatother
% Allow footnotes in longtable head/foot
\IfFileExists{footnotehyper.sty}{\usepackage{footnotehyper}}{\usepackage{footnote}}
\makesavenoteenv{longtable}
\setlength{\emergencystretch}{3em} % prevent overfull lines
\providecommand{\tightlist}{%
  \setlength{\itemsep}{0pt}\setlength{\parskip}{0pt}}
\setcounter{secnumdepth}{-\maxdimen} % remove section numbering
\ifLuaTeX
\usepackage[bidi=basic]{babel}
\else
\usepackage[bidi=default]{babel}
\fi
\babelprovide[main,import]{english}
\ifPDFTeX
\else
\babelfont{rm}[]{Latin Modern Roman}
\fi
% get rid of language-specific shorthands (see #6817):
\let\LanguageShortHands\languageshorthands
\def\languageshorthands#1{}
\usepackage{amsmath,amssymb}
\setlength{\emergencystretch}{2em}
\raggedbottom
\clubpenalty=10000
\widowpenalty=10000
\displaywidowpenalty=10000
\AtBeginDocument{\hypersetup{pdfauthor={},pdfsubject={},pdfkeywords={}}}
\makeatletter
\renewcommand{\maketitle}{%
  \begin{center}
    {\LARGE\@title\par}
    \vspace{0.8em}
    {\normalsize\@date\par}
  \end{center}
  \vspace{0.8em}
}
\makeatother

\predisplaypenalty=10000
\postdisplaypenalty=0
\interdisplaylinepenalty=10000
\ifLuaTeX
  \usepackage{selnolig}  % disable illegal ligatures
\fi
\usepackage{bookmark}
\IfFileExists{xurl.sty}{\usepackage{xurl}}{} % add URL line breaks if available
\urlstyle{same}
\hypersetup{
  pdftitle={Real rooted edge quotients and clique attachment for theta graphs},
  pdflang={en},
  hidelinks,
  pdfcreator={LaTeX via pandoc}}

\title{Real rooted edge quotients and clique attachment for theta
graphs}
\author{}
\date{2 November 2026}

\begin{document}
\maketitle

\subsection{Abstract}\label{abstract}

For every simple theta graph whose three path lengths are at least two,
we prove that the Bell transform of its unsigned chromatic edge quotient
has only strictly negative real zeros. Consequently, ordinary attachment
of any such theta graph to an arbitrary \(\tau\)-real graph along a
complete intersection preserves \(\tau\)-reality. The proof is exact and
computer-assisted. A rational radius criterion, together with two
Bernstein-certified factor-repair lemmas, covers every total path length
at least \(72\) and many smaller triples. An exhaustive finite
verification covers the remainder: among \(52,394\) sorted triples with
maximum path length at most \(68\), \(49,934\) satisfy the radius
criterion and \(2,460\) have exact negative-root certificates. The
finite verification checks \(83,842\) factor-root boxes using integer
arithmetic and rational sign changes, not numerical root classification.

This is a theta-family result, not a proof of general clique-overlap
closure for arbitrary \(\tau\)-real graphs. No historical novelty, human
peer review, or proof-assistant verification is asserted.

\subsection{1 Statement and
normalization}\label{statement-and-normalization}

All graphs are finite and simple. For \(a,b,c\geq2\), let
\(\Theta(a,b,c)\) consist of three internally vertex-disjoint paths of
lengths \(a,b,c\) with common distinct endpoints. Path lengths count
edges. For a graph \(G\) on \(n\) vertices, set
\nopagebreak
\[
P_G(x)=(-1)^n\chi_G(-x),
\]

where \(\chi_G\) is the chromatic polynomial. Define the linear Bell
transform by
\nopagebreak
\[
B((x)_j)=z^j,\qquad
(x)_j=x(x-1)\cdots(x-j+1),\qquad (x)_0=1.
\]

Equivalently,
\nopagebreak
\[
B(x^j)=T_j(z)=\sum_{i=0}^j S(j,i)z^i,
\]

where \(S(j,i)\) are Stirling numbers of the second kind and \(T_j\) are
the Touchard polynomials.

\textbf{Theorem.} For all integers \(a,b,c\geq2\), with
\(H=\Theta(a,b,c)\), the polynomial
\nopagebreak
\[
B(Q_H),\qquad Q_H(x)=\frac{P_H(x)}{x(x+1)},
\]

has only strictly negative real zeros. If \(G\) is any ordinary
\(\tau\)-real graph and \(G\cap H\) is complete, then the ordinary union
\(G\cup H\) is \(\tau\)-real.

Ordinary attachment retains the common clique, makes no additional
identifications, and adds no cross edges between the exclusive sides.
Since \(H\) is triangle-free, the possible complete intersections are
empty, one vertex, or one edge.

\subsubsection{1.1 The BRW transforms}\label{the-brw-transforms}

In Brenti--Royle--Wagner {[}1{]}, the transforms \(S\) and \(T\) send
falling and rising factorials, respectively, to monomials:
\nopagebreak
\[
S((x)_j)=z^j,\qquad
T\bigl(x(x+1)\cdots(x+j-1)\bigr)=z^j.
\]

Thus \(B=S\). If \(\varepsilon\) denotes reflection,
\((\varepsilon f)(x)=f(-x)\), then \(B=\varepsilon T\varepsilon\), and
\nopagebreak
\[
B(P_G)(z)=(-1)^n\tau_{\mathrm{BRW}}(G,-z).\tag{1}
\]

This is the reflection relation of {[}1, Proposition 5.1(d){]}.
Therefore real-rootedness of \(B(P_G)\) is exactly ordinary BRW
\(\tau\)-reality, not \(\sigma\)-reality of the original chromatic
polynomial. When these roots are real, they are nonpositive:
deletion--contraction gives \(P_G=P_{G-e}+P_{G/e}\), so induction shows
that \(P_G\), and hence \(B(P_G)\), has nonnegative coefficients.

We use the Bell product
\nopagebreak
\[
(f*g)(z)=\sum_{j\geq0}\frac{z^j}{j!}f^{(j)}(z)g^{(j)}(z).
\]

By {[}1, Propositions 4.2 and 4.4{]},
\nopagebreak
\[
B(FG)=B(F)*B(G),\tag{2}
\]

and \(f*g\) has only real nonpositive zeros whenever \(f\) and \(g\) do.
These are the root-preservation inputs used throughout the proof.

\subsection{2 The exact theta quotient and its root
trace}\label{the-exact-theta-quotient-and-its-root-trace}

Write
\nopagebreak
\[
S_\ell(t)=1+t+\cdots+t^{\ell-1}\quad(\ell\geq1),
\qquad S_0(t)=0,\qquad t=x+1.
\]

For a path of length \(\ell\), the numbers of proper internal colorings
with prescribed equal or prescribed distinct endpoint colors are
\nopagebreak
\[
E_\ell(q)=\frac{(q-1)^\ell+(q-1)(-1)^\ell}{q},\qquad
D_\ell(q)=\frac{(q-1)^\ell-(-1)^\ell}{q}.
\]

For example, these follow from
\nopagebreak
\[
E_{\ell+1}=(q-1)D_\ell,\qquad
D_{\ell+1}=E_\ell+(q-2)D_\ell,
\qquad E_0=1,\quad D_0=0.
\]

The displayed quotients are polynomials after cancellation. Choosing the
two common endpoint colors gives
\nopagebreak
\[
\chi_H(q)=qE_a(q)E_b(q)E_c(q)
+q(q-1)D_a(q)D_b(q)D_c(q).
\]

Put \(s=a+b+c\), so \(|V(H)|=s-1\). Substituting \(q=1-t=-x\) and
applying the sign \((-1)^{s-1}\) gives
\nopagebreak
\[
P_H(x)=\frac{t(t^a-1)(t^b-1)(t^c-1)
-(t^a-t)(t^b-t)(t^c-t)}{(t-1)^2}.
\]

After division by \(x(x+1)=(t-1)t\), the exact quotient identities are
\nopagebreak
\[
Q_H(t-1)=S_{a+b-1}(t)S_c(t)
+tS_{a-1}(t)S_{b-1}(t),\tag{3}
\]

and, equivalently,
\nopagebreak
\[
(t-1)^2Q_H(t-1)
=t^{s-1}-t^a-t^b-t^c+t+1.\tag{4}
\]

These are polynomial identities, including at \(t=1\). Formula (3) is
positive for every real \(t\geq0\), including \(t=0\), where its value
is \(1\). Thus every real \(t\)-root is negative. After \(t=x+1\), all
coefficients of \(Q_H(x)\) are strictly positive: the leading monomial
in (3) alone contributes positively to every coefficient.

The degree of \(Q_H(t-1)\) is
\nopagebreak
\[
d=s-3.
\]

It is monic and its next coefficient is \(2\). To see this, write
\nopagebreak
\[
N(t)=S_{a+b-1}(t)t^c-\bigl(S_a(t)+S_b(t)-1\bigr).
\]

Then \(Q_H(t-1)=N(t)/(t-1)\). The leading two coefficients of \(N\) are
\(1,1\), since \(a,b,c\geq2\) place the subtracted polynomial below
those terms. Division by \(t-1\) gives the asserted next coefficient.
Hence Vieta's formula gives
\nopagebreak
\[
\sum_{Q_H(t-1)=0}t=-2,\tag{5}
\]

with all roots counted with multiplicity.

\subsection{3 An effective radius
criterion}\label{an-effective-radius-criterion}

Set
\nopagebreak
\[
R=\frac{21}{20}.
\]

We first prove the analytic part of the theorem under the strict
rational condition
\nopagebreak
\[
R^{s-1}>R^a+R^b+R^c+R+1.\tag{6}
\]

If a root \(t\) of \(Q_H(t-1)\) had \(r=|t|\geq R\), equation (4) and
the triangle inequality would imply
\nopagebreak
\[
1\leq r^{a-s+1}+r^{b-s+1}+r^{c-s+1}
+r^{2-s}+r^{1-s}.
\]

Every exponent on the right is negative. Its value is therefore at most
its value at \(R\), which (6) makes strictly less than one. This
contradiction proves
\nopagebreak
\[
|t|<R\quad\text{for every root of }Q_H(t-1).\tag{7}
\]

\subsubsection{3.1 Every total path length at least 72 satisfies the
criterion}\label{every-total-path-length-at-least-72-satisfies-the-criterion}

For \(r>1\) and \(a,b\geq2\),
\nopagebreak
\[
r^a+r^b\leq r^{a+b-2}+r^2,
\]

because the difference is \(r^2(r^{a-2}-1)(r^{b-2}-1)\geq0\). Applying
this again gives
\nopagebreak
\[
r^a+r^b+r^c\leq r^{s-4}+2r^2.
\]

The exact integer inequality
\nopagebreak
\[
21^{68}(21^3-20^3)
>(2\cdot21^2+21\cdot20+20^2)20^{69}\tag{8}
\]

is equivalent to \(R^{68}(R^3-1)>2R^2+R+1\). Therefore, for \(s\geq72\),
\nopagebreak
\[
R^{s-1}>R^{s-4}+2R^2+R+1
\geq R^a+R^b+R^c+R+1.
\]

Thus every total path length at least \(72\) satisfies (6). This is a
concrete exact bound, not an asymptotic or existential threshold.

\subsection{4 Exact factor repair in the radius
disk}\label{exact-factor-repair-in-the-radius-disk}

Assume (7), and factor \(Q_H\) over \(\mathbb R\), keeping
multiplicities. A real \(t\)-root is negative by (3), so the
corresponding factor in \(x=t-1\) is \(x+\alpha\), with
\nopagebreak
\[
1<\alpha<1+R=\frac{41}{20}.
\]

A nonreal conjugate \(t\)-pair gives a monic factor
\nopagebreak
\[
h(x)=x^2+px+q,\qquad
p=2-2\operatorname{Re}t,\qquad q=|t-1|^2>0.
\]

Its Bell transform is \(B(h)=z^2+(p+1)z+q\). Call the pair \emph{good}
if
\nopagebreak
\[
J=(p+1)^2-4q\geq0,
\]

and \emph{bad} otherwise. For a good pair,
\(p+1=3-2\operatorname{Re}t>3-2R=9/10>0\), so \(B(h)\) has real strictly
negative roots.

For a bad pair, set
\nopagebreak
\[
\delta=q-p=|t|^2-1.
\]

Then
\nopagebreak
\[
J=(p-1)^2-4\delta<0,
\qquad
0<\delta<R^2-1=\frac{41}{400}<\frac19.
\]

Consequently \(|p-1|<2\sqrt\delta<2/3\) and
\nopagebreak
\[
\operatorname{Re}t=1-p/2>1/6.\tag{9}
\]

The next two lemmas repair products containing bad factors. Their
discriminant claims are certified by finite exact rational arithmetic
over entire parameter boxes.

\subsubsection{4.1 The tensor Bernstein certificate
principle}\label{the-tensor-bernstein-certificate-principle}

For a variable \(y\in[0,1]\), the degree-\(n\) Bernstein polynomials are
\nopagebreak
\[
b_{i,n}(y)=\binom ni y^i(1-y)^{n-i}\qquad(0\leq i\leq n).
\]

They are nonnegative and sum to one. For a multivariate polynomial of
coordinate degrees \(\mathbf n=(n_1,\ldots,n_r)\),
\nopagebreak
\[
D(\mathbf y)=\sum_{\mathbf j}c_{\mathbf j}\mathbf y^{\mathbf j}
=\sum_{\mathbf i}\beta_{\mathbf i}
\prod_{\ell=1}^r b_{i_\ell,n_\ell}(y_\ell),
\]

the exact conversion is
\nopagebreak
\[
\beta_{\mathbf i}
=\sum_{\mathbf j\leq\mathbf i}c_{\mathbf j}
\prod_{\ell=1}^r
\frac{\binom{i_\ell}{j_\ell}}{\binom{n_\ell}{j_\ell}}.\tag{10}
\]

It follows by applying, in each coordinate,
\nopagebreak
\[
y^j=\sum_{i=j}^n
\frac{\binom ij}{\binom nj}\binom ni y^i(1-y)^{n-i}.
\]

The coefficients can be checked by the inverse identity
\nopagebreak
\[
c_{\mathbf j}
=\left(\prod_{\ell=1}^r\binom{n_\ell}{j_\ell}\right)
\sum_{\mathbf i\leq\mathbf j}
(-1)^{|\mathbf j-\mathbf i|}\beta_{\mathbf i}
\prod_{\ell=1}^r\binom{j_\ell}{i_\ell},\tag{11}
\]

which is obtained by expanding the Bernstein basis back into powers. All
inequalities between multi-indices are coordinatewise.

If every \(\beta_{\mathbf i}>0\), then \(D\) is strictly positive
throughout \([0,1]^r\). For the two discriminants below, the supplied
certificates give every rational Bernstein coefficient, verify its
positivity, and verify the exact inverse conversion. Positivity of a
quartic discriminant alone does not imply four real roots. We also use a
real-rooted basepoint and connectedness: in a continuous monic
real-polynomial family with no repeated roots, the number of real roots
cannot change. The parameter boxes are connected and the discriminants
do not vanish.

\subsubsection{4.2 A bad quadratic times any conjugate
quadratic}\label{a-bad-quadratic-times-any-conjugate-quadratic}

Let \(h_1=x^2+Px+Q\) be a bad factor arising from the radius-\(R\) disk,
and let \(h_2=x^2+px+q\) be any conjugate-pair factor from that disk.
Then \(B(h_1h_2)\) has four distinct strictly negative roots.

To certify this, enlarge the parameter domain to \(a_0,b_0,u,v\in[0,1]\)
and set
\nopagebreak
\[
\begin{aligned}
P&=1+\frac23b_0(2a_0-1),&
Q&=P+\frac{b_0^2}{9},\\
p&=2+2R-4Ru,&
q&=(1+R-2Ru)^2+4R^2u(1-u)v.
\end{aligned}\tag{12}
\]

For the desired bad factor, take \(b_0=3\sqrt\delta\) and
\(2a_0-1=(P-1)/(2\sqrt\delta)\). The inequalities preceding (9) put
these parameters in the box. For the second root, use
\nopagebreak
\[
\operatorname{Re}t=-R+2Ru,\qquad
(\operatorname{Im}t)^2=4R^2u(1-u)v.
\]

This represents the entire disk. The added boundary cases only enlarge
the domain.

Direct Bell expansion gives
\nopagebreak
\[
\begin{aligned}
f(z)=B(h_1h_2)
={}&z^4+(6+P+p)z^3\\
&+(7+3(P+p)+Q+q+Pp)z^2\\
&+(1+P+p+Q+q+Pp+Pq+pQ)z+Qq.
\end{aligned}\tag{13}
\]

After substitution (12), \(\operatorname{Disc}_z(f)\) has coordinate
degrees
\nopagebreak
\[
(10,5,6,10)\quad\text{in }(u,v,a_0,b_0).
\]

All \(11\cdot6\cdot7\cdot11=5,082\) tensor Bernstein coefficients are
strictly positive. Their exact minimum is
\nopagebreak
\[
\frac{1319869198532766581}{19683000000000000}>0.\tag{14}
\]

The complete rational list is supplied in
\texttt{\nolinkurl{effective_mixed_certificate.json}};
\texttt{\nolinkurl{effective_mixed_certificate.py}} constructs the
discriminant, performs (10), and checks the exact inverse expansion
(11).

At \(b_0=0\), \(u=31/42\), and \(v=300/341\), independently of \(a_0\),
both factors are \(x^2+x+1\), and
\nopagebreak
\[
\begin{aligned}
f(z)&=z^4+8z^3+16z^2+8z+1\\
&=\bigl(z^2+(4+\sqrt2)z+1\bigr)
\bigl(z^2+(4-\sqrt2)z+1\bigr).
\end{aligned}\tag{15}
\]

This has four distinct strictly negative roots. Nonvanishing
discriminant and connectedness therefore give four distinct real roots
throughout the parameter box.

Moreover, \(P\geq1/3\), \(Q\geq P\), \(p\geq-1/10\), and \(q\geq0\).
These inequalities make every nonconstant coefficient in (13) strictly
positive and the constant nonnegative. For example, the coefficient of
\(z\) is
\nopagebreak
\[
1+p+P(1+p)+Q(1+p)+Pq+q
\geq\frac9{10}+\frac9{10}P+\frac9{10}Q>0.
\]

Thus no root is positive. In the actual nonreal-pair application,
\(q,Q>0\), so no root is zero either. This proves the lemma. On the
enlarged box boundary \(q=0\), a simple zero is allowed; strict
negativity is asserted for the actual factors used in the proof.

\subsubsection{4.3 A bad quadratic times a real linear
factor}\label{a-bad-quadratic-times-a-real-linear-factor}

Let
\nopagebreak
\[
h(x)=x^2+Ax+(A+e),\qquad
0\leq e\leq\frac19,\qquad (A-1)^2\leq4e,
\]

and let \(1\leq\alpha\leq21/10\). Then \(B(h(x)(x+\alpha))\) has three
distinct strictly negative roots.

The transformed cubic is
\nopagebreak
\[
f(z)=z^3+(3+A+\alpha)z^2
+(1+2A+\alpha+\alpha A+e)z+\alpha(A+e).\tag{16}
\]

Parameterize the region by \(u,v,w\in[0,1]\):
\nopagebreak
\[
y=2u-1,\qquad
A=1+\frac23y,\qquad
e=\frac{y^2+(1-y^2)v}{9},\qquad
\alpha=1+\frac{11}{10}w.\tag{17}
\]

This covers the stated curved \((A,e)\) region; when \(|y|=1\),
necessarily \(e=1/9\), and \(v\) can be chosen arbitrarily. The exact
discriminant in \((u,v,w)\) has coordinate degrees \((6,3,4)\). Its
\(7\cdot4\cdot5=140\) Bernstein coefficients are all positive, with
exact minimum
\nopagebreak
\[
\frac{72882116}{6834375}>0.\tag{18}
\]

The full list and the exact inverse conversion are supplied in
\texttt{\nolinkurl{effective_cubic_certificate.json}} and checked by
\texttt{\nolinkurl{effective_cubic_certificate.py}}. At
\((u,v,w)=(1/2,0,0)\), one has \(A=1\), \(e=0\), \(\alpha=1\), and
\nopagebreak
\[
f(z)=(z+1)(z^2+4z+1),
\]

with three distinct negative roots. Connectedness and the nonvanishing
discriminant give three distinct real roots everywhere in the cube. All
coefficients are strictly positive because \(A\geq1/3\), \(e\geq0\), and
\(\alpha\geq1\), so all roots are strictly negative.

Every actual bad quadratic from (7) satisfies these conditions with
\(A=p\) and \(e=\delta\). Every actual real linear factor also fits
because \(1<\alpha<41/20<21/10\).

\subsection{5 Assembly for every radius certified
triple}\label{assembly-for-every-radius-certified-triple}

If the number of bad quadratic factors is even, pair them arbitrarily
and use Section 4.2 on each pair. If it is odd and a real linear factor
exists, combine one bad quadratic with that real factor using Section
4.3, then pair the remaining bad quadratics.

Suppose instead that the number is odd and there are no real factors.
There must be at least one good quadratic: otherwise all \(t\)-roots
would have positive real part by (9), contradicting their total sum
\(-2\) in (5). Reserve one good quadratic and combine it with one bad
quadratic using Section 4.2, then pair the remaining bad quadratics.

Unused good quadratics and real linears already have strictly negative
Bell roots. Repeated factors are counted with multiplicity throughout;
each factor occurrence is used once. Iterating (2) proves that
\(B(Q_H)\) has only real nonpositive zeros. Its constant coefficient is
\(Q_H(0)>0\), by (3) at \(t=1\), so zero is not a root. This proves the
theorem for every triple satisfying (6), including every total path
length at least \(72\).

\subsection{6 Exhaustive exact verification of the finite
remainder}\label{exhaustive-exact-verification-of-the-finite-remainder}

\subsubsection{6.1 The finite domain and its exact
partition}\label{the-finite-domain-and-its-exact-partition}

By symmetry, sort the path lengths as \(2\leq a\leq b\leq c\). If
\(c\geq69\), then \(a+b+c\geq73\), so Section 3.1 applies. It therefore
suffices to enumerate
\nopagebreak
\[
2\leq a\leq b\leq c\leq68.
\]

There are \(\binom{69}{3}=52,394\) such triples. For each triple, put
\(n=a+b+c-1\). Clearing positive denominators in (6) gives the integer
margin
\nopagebreak
\[
\begin{aligned}
\Delta(a,b,c)={}&21^n
-21^a20^{n-a}-21^b20^{n-b}-21^c20^{n-c}\\
&-21\cdot20^{n-1}-20^n.
\end{aligned}\tag{19}
\]

A positive margin selects the analytic branch. Every other triple must
have a complete exact root certificate. The exhaustive loop uses
\(c=2,\ldots,68\), \(b=2,\ldots,c\), and \(a=2,\ldots,b\), so no
ordering or symmetry case is omitted.

\subsubsection{6.2 Independent integer constructions of the intended
polynomial}\label{independent-integer-constructions-of-the-intended-polynomial}

Write
\nopagebreak
\[
Q_H(x)=\sum_{j=0}^d c_jx^j,\qquad
B(Q_H)(z)=\sum_{i=0}^d b_iz^i,
\qquad d=a+b+c-3.
\]

One direct coefficient formula follows from (4), after \(t=x+1\) and
division by \(x^2\):
\nopagebreak
\[
c_j=\binom{a+b+c-1}{j+2}
-\binom a{j+2}-\binom b{j+2}-\binom c{j+2}.\tag{20}
\]

Binomial coefficients outside their usual range are zero. The linear and
constant terms \(t+1\) in (4) contribute nothing to these coefficients.

A separate construction starts with the two geometric-sum products in
(3). If their convolution gives
\nopagebreak
\[
Q_H(t-1)=\sum_{k=0}^d q_kt^k,
\]

then the shift and Bell transform are
\nopagebreak
\[
c_j=\sum_{k=j}^d q_k\binom kj,\qquad
b_i=\sum_{j=i}^d c_jS(j,i).\tag{21}
\]

All these operations use integers. The Stirling numbers can be generated
from \(S(0,0)=1\) and
\nopagebreak
\[
S(j+1,i)=iS(j,i)+S(j,i-1).
\]

The independent finite auditor reconstructs the intended polynomial by
the geometric-convolution route, checks its degree \(d\), monicity and
positive coefficients, and compares it with the certificate's polynomial
identity. This is not a verification that merely accepts polynomial
coefficients supplied by a root-finding program.

\subsubsection{6.3 The exact negative root certificate
criterion}\label{the-exact-negative-root-certificate-criterion}

For each remaining triple, the certificate lists monic integer
polynomials \(f_\nu(z)\) and positive integer multiplicities \(m_\nu\).
The checker verifies the exact coefficient identity
\nopagebreak
\[
B(Q_H)(z)=\prod_\nu f_\nu(z)^{m_\nu}.\tag{22}
\]

For each listed factor of degree \(e\), it requires exactly \(e\)
pairwise strictly disjoint rational boxes, ordered along the negative
real axis. A box is either:

\begin{itemize}
\tightlist
\item
  a rational singleton \(r<0\) with \(f_\nu(r)=0\) exactly; or
\item
  an interval \([\ell,u]\) with \(\ell<u<0\) and
  \(f_\nu(\ell)f_\nu(u)<0\) exactly.
\end{itemize}

Strict disjointness is required within each factor. The intermediate
value theorem supplies a distinct negative root in each box. Since there
are degree-many boxes, these exhaust all roots of that factor and make
them simple. Boxes belonging to different factors need not be disjoint:
the exact product identity (22) and the recorded multiplicities handle
repeated roots of the full polynomial.

Endpoint signs use homogeneous integer evaluation. For
\(f(z)=\sum_{j=0}^e f_jz^j\) and a rational endpoint \(n/d\) with
\(d>0\),
\nopagebreak
\[
\mathcal E_f(n,d)=\sum_{j=0}^e f_jn^jd^{e-j}
=d^ef(n/d).\tag{23}
\]

Thus the integer \(\mathcal E_f(n,d)\) has exactly the sign of
\(f(n/d)\). The checker uses this identity, exact integer polynomial
multiplication, and degree counting. It does not trust a factorization
oracle, a numerical real/nonreal classification, or an approximate
root-isolator conclusion. Numerical or interval software may propose
candidate boxes during certificate generation; a successful verification
depends only on the exact tests just stated.

\subsubsection{6.4 The completed finite
verification}\label{the-completed-finite-verification}

The supplied finite certificates and coverage record pass the complete
check with the following results:

\begin{longtable}[]{@{}lr@{}}
\toprule\noalign{}
Quantity & Exact count \\
\midrule\noalign{}
\endhead
\bottomrule\noalign{}
\endlastfoot
Sorted triples in the finite domain & 52,394 \\
Triples satisfying the strict radius test & 49,934 \\
Triples with complete exact root certificates & 2,460 \\
Missing triples & 0 \\
Factor-root boxes checked & 83,842 \\
\end{longtable}

The \(83,842\) count is the sum of the degrees of the listed factors,
without multiplying by their recorded exponents; it is not a count of
globally distinct numerical roots across different triples. Every
exceptional quotient has degree at most \(68\), and the largest
exceptional path length is \(67\).

\begin{samepage}
The complete records are supplied in
\texttt{\nolinkurl{finite_certificates.zip}}, with file identities in
\texttt{\nolinkurl{certificate_manifest.json}} and one coverage entry
per finite-domain triple in
\texttt{\nolinkurl{finite_coverage.jsonl.gz}}. The replay recomputes
(19) for every triple. It rejects missing certificates, certificates
assigned to the wrong triple, failed factor or root-box identities, and
unexpected certificate files. The counts are consistency checks, not
substitutes for exhaustive coverage.

\end{samepage}

It follows from (22)--(23) that every one of the \(2,460\) remaining
Bell quotients has only strictly negative real zeros. Together with the
\(49,934\) analytic branches and the unbounded analytic tail of Section
3.1, this exhausts all \(a,b,c\geq2\) and proves the quotient assertion
of the theorem.

\subsection{7 Clique attachment}\label{clique-attachment}

If \(G\cap H=K_2\), chromatic complete-cutset factorization {[}1,
Proposition 1.1{]} gives
\nopagebreak
\[
P_{G\cup H}=P_G\frac{P_H}{x(x+1)}=P_GQ_H.
\]

For an ordinary \(\tau\)-real \(G\), formula (1) makes \(B(P_G)\)
real-rooted with nonpositive zeros. The proved quotient theorem and (2)
then show that \(G\cup H\) is \(\tau\)-real.

The theta graph itself is also \(\tau\)-real, since
\nopagebreak
\[
P_H=x(x+1)Q_H,\qquad B(x(x+1))=z(z+2),
\]

and the Bell-product theorem applies once more. BRW Proposition 5.12
therefore gives the empty and single-vertex overlap cases. Every cycle
of \(H\) uses two of its paths and has length at least four, so \(H\) is
triangle-free. These cases exhaust its possible complete overlaps and
finish the theorem.

\subsection{8 Reproducibility and
limitations}\label{reproducibility-and-limitations}

This is an exact computer-assisted proof. The two repair lemmas use
complete rational Bernstein certificates for explicit discriminants over
compact parameter boxes. The finite part uses exact rational sign
changes and integer factor identities. Both computational components are
proof obligations, not merely illustrative numerical checks.

The accompanying \texttt{\nolinkurl{verify_all.py}},
\texttt{\nolinkurl{check_certificates.py}}, and
\texttt{\nolinkurl{audit_finite_certificates.py}} provide the package
verification route. The two effective-certificate scripts reconstruct
the Bernstein calculations; \texttt{\nolinkurl{finite_theta_verify.py}}
supplies the finite certificate generator. The mathematical acceptance
conditions are stated explicitly in Sections 4.1 and 6.3, so the claims
are not based on the generator's success message alone. Detailed
execution and integrity information is given in the accompanying package
documentation.

AI systems assisted with derivation, review passes, exact computations
and manuscript preparation. AI review is not independent human peer
review, and this work is not proof-assistant formalization. The theorem
concerns the stated theta family; it does not establish the full
\(\tau\) clause of BRW Conjecture 6.5 for arbitrary pairs of
\(\tau\)-real graphs.

\subsubsection{8.1 An unresolved literature
comparison}\label{an-unresolved-literature-comparison}

A potentially relevant paper is Lili Mu, Qi Shen and Xinyao Xu, \emph{On
the Real-rootedness, invariance, and root bounds of
\(\tau\)-polynomials}, Discrete Mathematics, Algorithms and
Applications, article 2550185 {[}2{]}. Its publisher-deposited Crossref
record lists an online-publication date of 9 January 2026; some
secondary listings show a 2025 date.

\begin{samepage}
The full text, exact normalization, graph families, hypotheses and
theorem statements of that paper have not been inspected. The accessible
deposited abstract has mathematical expressions replaced by formula
placeholders, and the publisher text was not accessible in the bounded
literature check. Whether it subsumes, overlaps with, or differs from
the present quotient and attachment theorem therefore remains unknown.
Inaccessible formulas are not evidence of non-overlap. No historical
priority or comprehensive-literature certificate is claimed.

\end{samepage}

\subsection{References}\label{references}

\begin{enumerate}
\def\labelenumi{\arabic{enumi}.}
\item
  Francesco Brenti, Gordon F. Royle, and David G. Wagner. \emph{Location
  of zeros of chromatic and related polynomials of graphs}. Canadian
  Journal of Mathematics \textbf{46} (1), 55--80 (1994).
  \href{https://doi.org/10.4153/CJM-1994-002-3}{doi:10.4153/CJM-1994-002-3}.
  Proposition 1.1 gives complete-cutset chromatic factorization;
  Propositions 4.2 and 4.4 give the Bell product and its root
  preservation; Proposition 5.1(d) gives reflection; Proposition 5.12
  covers intersections of at most one vertex.
\item
  Lili Mu, Qi Shen, and Xinyao Xu. \emph{On the Real-rootedness,
  invariance, and root bounds of \(\tau\)-polynomials}. Discrete
  Mathematics, Algorithms and Applications, article 2550185.
  \href{https://doi.org/10.1142/S179383092550185X}{doi:10.1142/S179383092550185X}.
  Bibliographic metadata:
  \href{https://api.crossref.org/works/10.1142/S179383092550185X}{publisher-deposited
  Crossref record}. The precise mathematical scope remains unverified,
  as explained in Section 8.1.
\end{enumerate}

\end{document}
